\documentclass[12pt,a4paper,oneside]{book}

% =====================================================================
% PACOTES ESSENCIAIS
% =====================================================================
\usepackage[brazilian]{babel}
\usepackage[T1]{fontenc}
\usepackage[utf8]{inputenc}
\usepackage{amsmath}
\usepackage{amssymb}
\usepackage{amsthm}
\usepackage{geometry}
\usepackage{fancyhdr}
\usepackage{hyperref}
\usepackage{booktabs}
\usepackage{mathtools}

% =====================================================================
% GEOMETRIA
% =====================================================================
\geometry{
    left=2.5cm,
    right=2.5cm,
    top=3cm,
    bottom=3cm,
    headheight=14pt
}

% =====================================================================
% CABECALHO E RODAPE
% =====================================================================
\pagestyle{fancy}
\fancyhf{}
\fancyhead[L]{\textit{Capitulo 01: Fundamentos de Modelagem Matematica}}
\fancyhead[R]{\thepage}
\fancyfoot[C]{\small Principios de Modelagem Matematica}
\renewcommand{\headrulewidth}{0.4pt}
\renewcommand{\footrulewidth}{0.4pt}

% =====================================================================
% AMBIENTES
% =====================================================================
\theoremstyle{definition}
\newtheorem{definition}{Definicao}[chapter]
\newtheorem{example}[definition]{Exemplo}
\newtheorem{remark}[definition]{Observacao}
\newtheorem{problem}{Problema}

% =====================================================================
% DOCUMENTO
% =====================================================================
\title{\textbf{Capitulo 01: Fundamentos de Modelagem Matematica}\\[0.4cm]
\large Principios de Modelagem Matematica}
\author{Baseado em Dym, C. L.}
\date{2026-04-14}

\begin{document}

\maketitle
\tableofcontents
\newpage

\chapter{Fundamentos de Modelagem Matematica}

\section{Ideia central do capitulo}

Modelagem matematica e o processo de representar um sistema real por meio de variaveis,
hipoteses e equacoes, com objetivo de explicar, prever ou otimizar comportamentos.

Em termos simples, modelar e traduzir uma pergunta do mundo real para a linguagem da matematica.

\section{O que e um modelo?}

\begin{definition}[Modelo matematico]
Um modelo matematico e uma representacao simplificada de um fenomeno real,
construida com variaveis, parametros, relacoes funcionais e condicoes de contorno/iniciais.
\end{definition}

Todo modelo envolve compromisso entre realismo e simplicidade.

\begin{remark}
Um modelo nao precisa ser perfeito para ser util. Ele precisa ser adequado ao objetivo.
\end{remark}

\section{Etapas classicas da modelagem}

\begin{enumerate}
    \item Definir o problema e o objetivo da analise.
    \item Identificar variaveis relevantes e parametros.
    \item Formular hipoteses simplificadoras.
    \item Escrever equacoes do modelo.
    \item Resolver analitica ou numericamente.
    \item Validar com dados reais.
    \item Refinar o modelo, se necessario.
\end{enumerate}

\subsection{Exemplo de fluxo pratico}

Problema: prever a temperatura de um cafe ao longo do tempo.

\begin{itemize}
    \item Variavel de estado: $T(t)$.
    \item Parametros: temperatura ambiente $T_a$, constante de troca termica $k$.
    \item Hipotese: taxa de resfriamento proporcional a diferenca $T-T_a$.
\end{itemize}

Modelo:
\begin{equation}
\frac{dT}{dt}=-k(T-T_a).
\end{equation}

Esse exemplo mostra como uma frase fisica vira equacao diferencial.

\section{Tipos de modelos}

\subsection{Por natureza das variaveis}

\begin{itemize}
    \item \textbf{Continuos}: variaveis definidas em intervalos reais (tempo continuo).
    \item \textbf{Discretos}: variaveis em passos (dias, iteracoes, ciclos).
\end{itemize}

\subsection{Por estrutura matematica}

\begin{itemize}
    \item \textbf{Algebricos}: sem derivadas.
    \item \textbf{Diferenciais ordinarios (EDO)}: derivadas no tempo.
    \item \textbf{Diferenciais parciais (EDP)}: derivadas no tempo e no espaco.
    \item \textbf{Estocasticos}: incluem aleatoriedade.
\end{itemize}

\subsection{Por objetivo}

\begin{itemize}
    \item \textbf{Descritivo}: explica o fenomeno.
    \item \textbf{Preditivo}: antecipa comportamento futuro.
    \item \textbf{Otimizacao}: busca melhor decisao sob restricoes.
\end{itemize}

\section{Hipoteses e parametros}

Uma boa modelagem depende de hipoteses claras.

Exemplos de hipoteses comuns:
\begin{itemize}
    \item sistema fechado (sem troca de massa);
    \item propriedades constantes no tempo;
    \item efeitos secundarios despreziveis.
\end{itemize}

\begin{definition}[Parametro]
Parametro e uma quantidade fixa dentro de uma simulacao, mas que pode mudar entre experimentos.
\end{definition}

\begin{definition}[Variavel de estado]
Variavel de estado descreve a condicao instantanea do sistema e evolui no tempo ou no espaco.
\end{definition}

\section{Exemplos introdutorios}

\subsection{Exemplo 1: crescimento populacional simples}

Hipotese: taxa de crescimento proporcional a populacao atual.

\begin{equation}
\frac{dP}{dt}=rP, \qquad P(0)=P_0.
\end{equation}

Solucao:
\begin{equation}
P(t)=P_0e^{rt}.
\end{equation}

Interpretacao:
\begin{itemize}
    \item $r>0$: crescimento.
    \item $r<0$: decaimento.
\end{itemize}

\subsection{Exemplo 2: queda livre sem resistencia do ar}

Hipoteses:
\begin{itemize}
    \item aceleracao constante $g$;
    \item sem arrasto do ar.
\end{itemize}

Modelo:
\begin{equation}
\frac{dv}{dt}=g, \qquad \frac{dy}{dt}=v.
\end{equation}

Com $v(0)=0$ e $y(0)=h$:
\begin{equation}
v(t)=gt, \qquad y(t)=h-\frac{1}{2}gt^2.
\end{equation}

Esse modelo e simples e util, mas perde precisao em altas velocidades por ignorar arrasto.

\subsection{Exemplo 3: mistura em tanque}

Tanque bem misturado de volume constante $V$, entrada de concentracao $c_{in}$ e vazao $q$.

Se $c(t)$ e concentracao no tanque:
\begin{equation}
V\frac{dc}{dt}=q(c_{in}-c).
\end{equation}

Esse modelo aparece em engenharia quimica, ambiental e biologica.

\section{Validacao e limites do modelo}

Validar e comparar previsoes com observacoes.

Indicadores comuns:
\begin{itemize}
    \item erro medio absoluto;
    \item erro relativo;
    \item tendencia sistematica de superestimar/subestimar.
\end{itemize}

\begin{remark}
Quando o erro e grande, nao significa que a matematica esta errada; geralmente as hipoteses precisam ser revistas.
\end{remark}

\section{Boas praticas para modelar}

\begin{enumerate}
    \item Comece com modelo minimo viavel (simples e interpretavel).
    \item Nomeie variaveis e unidades com clareza.
    \item Registre hipoteses explicitamente.
    \item Verifique consistencia dimensional das equacoes.
    \item So adicione complexidade se houver ganho comprovado.
\end{enumerate}

\section{Exercicios curtos}

\begin{problem}
Escreva um modelo simples para saldo bancario com deposito mensal constante e juros proporcionais ao saldo atual.
\end{problem}

\begin{problem}
No modelo de queda livre sem arrasto, determine o tempo para um objeto atingir o solo, partindo de altura $h$ e $v(0)=0$.
\end{problem}

\begin{problem}
Para o modelo de mistura em tanque, encontre o estado estacionario da concentracao.
\end{problem}

\section{Resumo final}

O capitulo 01 estabelece a base da modelagem matematica:
\begin{itemize}
    \item formular perguntas de forma quantitativa;
    \item transformar fenomenos em equacoes;
    \item explicitar hipoteses e parametros;
    \item validar e refinar com dados.
\end{itemize}

Com essa base, os capitulos seguintes aprofundam ferramentas especificas como analise dimensional, escala, aproximacoes e dinamicas exponenciais.

\end{document}
